\documentclass[10pt,a4paper]{amsart}
\usepackage[margin=19mm]{geometry}
\usepackage{amsmath,amssymb,mathtools}
\usepackage[scr=rsfs,cal=euler]{mathalfa}
\usepackage{xfrac}
\usepackage[unicode,hidelinks,bookmarks=false]{hyperref}
\newcommand{\ie}{i.e.\ }
\newcommand{\cf}{cf.\ }
\newcommand{\resp}{respectively\ }
\newcommand{\Subsection}{\subsection}
\providecommand{\nobreakdash}{\nobreak\hbox{-}}
\setlength{\emergencystretch}{2em}
\multlinegap0pt
\usepackage{amssymb}
\usepackage{tikz}
\usepackage{bm}
\newtheorem{lem}{Lemma}
\newcommand{\C}{\mathbb{C}}
\title{$p$-adic properties of division polynomials and algebraic sigma functions}
\author{Yu Katagiri, Shinichi Kobayashi}
\date{Source: arXiv:2609.07024v1 (CC BY 4.0)}
\begin{document}
\maketitle

\noindent\emph{Source and licence.} $p$-adic properties of division polynomials and algebraic sigma functions. Version arXiv:2609.07024v1 (CC BY 4.0), distributed by arXiv under the Creative Commons Attribution 4.0 International licence, http://creativecommons.org/licenses/by/4.0/, as recorded in the arXiv licence metadata for this exact version and shown on the abstract page https://arxiv.org/abs/2609.07024v1. Full licence notice: \texttt{licenses/arxiv-2609.07024-CC-BY-4.0.txt}.\par\smallskip

\noindent\textit{Editorial scope.} Counted unit: the article's lem pWPT (source0.tex lines 694--716). The article's complete original proof of it is delivered with it (source0.tex lines 718--748, one \texttt{proof} environment of 31 lines). No mathematics is written, repaired or re-derived: the statement and the proof are the article's own lines, in the article's own notation, and no retained line is substituted or re-flowed.  No further source-local context is needed: every cross-reference of every retained line resolves inside the two delivered spans.  Nothing was omitted from any retained span, so the package carries no omission record.  No retained line contains a citation, so the wrapper carries no bibliography and no bibliography entry.  No retained line contains a figure environment, an image-inclusion command or drawing code, so no image is delivered and the package declares no asset dependency.  Editorial formatting only, in addition: an AMS \texttt{amsart} preamble that restates the article's own declarations and macros used by the retained lines, namely the environments \texttt{lem} on separate counters, together with the standard packages those lines need.  The retained environments print in this document as Lemma 1 in order of appearance, whereas the article prints the same items with the numbering of its own full document.  The complete native source of the article as distributed in the arXiv e-print for this exact version is bundled unchanged as \texttt{sources/209599/source0.tex}.
\begin{lem}\label{pWPT}
Let $E$ be an elliptic curve over $K$ with good reduction, let
$f \in K(E)$ be a rational function, and let $P \in E(K)$.
Suppose that $f$ has no poles on the residue disc
$V_P$, the set of points of $E(\mathbb{C}_p)$ whose 
reduction equals that of $P$. 

\begin{enumerate}
\item Let $\mathcal{E}/\mathcal{O}_K$ be the smooth model of $E$, and let
$t$ be a local parameter along the section over $\mathcal{O}_K$ defined by
$P$. 
Regarding $f$ as an element of $K((t))$ via the formal completion along
this section, that is, via the embedding $K(E) \hookrightarrow K((t))$,
we have $f \in \mathcal{O}_K[[t]] \otimes K$.
\item Suppose that the order $r$ of the reduction of $P$ is prime to $p$,
and let $T$ be the Teichm\"{u}ller lift of $P$. Then we have
%
\begin{align*}
\lim_{\substack {k \to \infty \\ p^k \equiv 1 \bmod r}} f(p^kP)=f(T).
\end{align*}
%
\end{enumerate}
\end{lem}

\begin{proof}
\noindent\textup{(1)} By translating by $P$, we may assume that $P=O$.
We note that $t$ identifies the residue disc $V_O$ with the open unit disc $\{ a \in \C_p \mid |a|<1\}$. 
We may write
%
\begin{align*}
f(t)=\frac{g(t)}{h(t)}
\end{align*}
%
for some $g(t), h(t) \in \mathcal{O}_K[[t]]$.
(Indeed, $f$ is a rational function in $x$ and $y$, and $t^2x(t), t^3y(t) \in  \mathcal{O}_K[[t]]$.)
Since $f$ has no pole on the residue disc, by the $p$-adic Weierstrass preparation theorem, 
we may take 
$h(t)\in \mathcal{O}_K[[t]]^{\times}$ up to a factor in $K^{\times}$, and the assertion follows. 

\medskip
\noindent\textup{(2)} Let $t$ be a local parameter along the section defined by $T$ on the smooth model.
Then the residue disc $V_T=V_P$ is identified with $T+V_O$.
Write $P=T+R$ and let $t_R$ be the coordinate of $R$ in $V_O$.
If $p^k \equiv 1 \bmod r$, then $p^kT=T$ since $T \in E(K)[r]$, and hence $p^kP=T+p^kR$. 
The point $p^kR$ corresponds to the coordinate 
$[p^k]t_R \in \mathfrak{m}_K^k$. 
By (1) applied to $T$ in place of $P$, we have  $f(t) \in \mathcal{O}_K[[t]] \otimes K$ and  
%
\begin{align*}
\lim_{\substack {k \to \infty \\ p^k \equiv 1 \bmod r}} f(p^kP)
=\lim_{\substack {k \to \infty \\ p^k \equiv 1 \bmod r}} f([p^k]t_R)
=f(0)=f(T).
\end{align*}
%
\end{proof}

\end{document}
